\documentclass[]{article}
\usepackage[utf8]{inputenc}
\usepackage{textcomp}

\usepackage{graphicx}
\usepackage{float}

\usepackage{times}
\usepackage{color}
\usepackage{xcolor}
\usepackage{natbib} 
\usepackage{enumerate}
\usepackage{caption}
%\usepackage[separator = {;}]{credits}
\usepackage{lscape}
\usepackage{subcaption}
\usepackage{wrapfig}
\usepackage{orcidlink}
\usepackage{amssymb}
\usepackage{amsthm}
\usepackage{enumitem}
\usepackage{hyperref}
\hypersetup{colorlinks=true,
	citecolor=cyan,
	pdfauthor=cyan,
	linkcolor=cyan}

%opening
\title{}
\author{}

\begin{document}

\maketitle

\begin{abstract}

\end{abstract}

\section{Requirement definition}
\subsection{Mission objectives}
Environmental Monitoring: Monitor environmental factors such as deforestation, coastline changes, pollution levels, and natural disasters to support environmental management and conservation efforts.
\subsection{Mission requirements:}
	\begin{itemize}
		\item Spatial Resolution: Specify the desired spatial resolution of the imagery, typically measured in meters per pixel. This determines the level of detail captured in the images and influences the selection of the imaging sensor and optics.
		\item  Spectral Bands: Define the required spectral bands to capture information across different parts of the electromagnetic spectrum. This may include visible, near-infrared, and shortwave infrared bands for vegetation health analysis and other applications.
		\item  Field of View (FOV): Determine the desired FOV to cover a specific area on the Earth's surface in a single image. The FOV impacts the selection and design of the optics system.
		\item  Image Quality: Define the desired image quality in terms of radiometric accuracy, signal-to-noise ratio, and distortion levels. This ensures reliable and accurate interpretation of the acquired imagery.
		\item  Revisit Time: Specify the desired revisit time, which refers to how often the satellite can capture imagery of a specific location. Revisit time depends on the satellite's orbital parameters and the desired coverage and monitoring frequency.
		\item Data Storage and Downlink: Determine the data storage capacity of the CubeSat and the downlink bandwidth for transmitting acquired imagery back to the ground station. Consider the limitations of the CubeSat's resources and the data processing capabilities on board.
		\item  Power Consumption: Define the power requirements for the camera system, including the imaging sensor, optics, electronics, and communication subsystems. Ensure the power consumption is within the capabilities of the CubeSat's power supply.
		\item Size, Weight, and Volume: Consider the size, weight, and volume constraints of the CubeSat platform. Ensure that the camera system can be integrated within
	\end{itemize}
\section{System architecture}
\section{Imaging sensor selection}
\section{Optical camera design}
\section{Mechanical design}
\section{Electronics and control}
\section{Testing and calibration}
\section{Integration and verification}
\section{Mission operation and data analysis}

\end{document}
